\section{Electromagnetic Entry, Feedback, and Observation}

\subsection{Physical tangent and frozen inverse problem}

Consider isotropic, nonmagnetic dielectric contrast \(\chi\) with a fixed acquisition geometry. For illumination \(p\), the contrast-source formulation is

\[
\begin{aligned}
j_p&=X(e_p^{\rm inc}+G_Dj_p),&
L&=I-XG_D,\\
L\,\delta j_p&=B_ps,&
B_ps&=DX[Q_Ts]e_p^{\rm tot}.
\end{aligned}\tag{1}
\]

Here \(G_D\) is the internal vector Green operator; \(X\) is the material response; \(e_p^{\rm tot}\) is the current total field; and \(Q_T\) expands a real material coordinate vector \(s\in\mathbb R^d\), including both parts of complex contrast. Thus \(d\) is the admissible real material dimension. Its realified expansion obeys \(Q_T^\top M_\chi Q_T=I\), with the volume-weighted material \(L^2\) metric \(M_\chi\). Current coordinates are also normalized to their declared discrete metric. With the receiver/scaling map \(S\),

\[
K_p=L^{-1}B_p,\qquad J_p=SK_p. \tag{2}
\]

The actual discrete-dipole approximation (DDA) uses contrast-dependent polarizability, rather than substituting \(X=\chi\) in its discrete equation. Its injection retains the polarizability derivative and the vector local field; Section 1 of the Supplement gives the exact implemented model. The operator identities use \(A=I-L\); \(A=XG_D\) expresses their contrast-source interpretation, and \(A=D_aG_{\rm off}\) is the implemented DDA counterpart.

Illuminations have separate currents and one shared material vector $s\in\mathbb R^d$. Each excitation $B_ps$ produces a current change $L^{-1}B_ps$ and measurement change $J_ps$. Stacking $[\operatorname{Re}(J_ps);\operatorname{Im}(J_ps)]$ over illuminations gives one real measurement vector. The normal equations are therefore real material-space equations, with $^*$ denoting the corresponding adjoint. The complex current-block identities lift through block-diagonal current operators and vertically stacked material injection (Section 1 of the Supplement).

At a frozen state, let \(r\) be predicted minus observed data, \(\Lambda\succ0\) the common regularization/damping curvature, and \(\ell\) its common linear prior term. Define the full quadratic

\[
\begin{aligned}
\Phi(s)&=\frac12\|r+Js\|^2\\
&\quad+\frac12s^*\Lambda s+\ell^*s\\
&=\frac12s^*Hs+h^*s+\frac12\|r\|^2,\\
H&=J^*J+\Lambda,& h&=J^*r+\ell.
\end{aligned}\tag{3}
\]

Thus \(s_*=-H^{-1}h\), and the squared \(H\)-energy step error is twice the full quadratic gap, providing an exact measure of local step quality. Subtracting the common constant \(\tfrac12\|r\|^2\) in stored relative objectives leaves all gaps and gains unchanged.

\subsection{Specified Galerkin enlargement}

Let \(V\) be an orthonormal current basis, \(C\) an orthonormal candidate bank with \(V^*C=0\), and assume \(V^*LV\) is invertible. The old Galerkin operators are

\[
\begin{aligned}
R_0&=V(V^*LV)^{-1}V^*,& K_0&=R_0B,\\
J_0&=SK_0,& H_0&=J_0^*J_0+\Lambda,\\
&&M_0&=H_0^{-1}.
\end{aligned}\tag{4}
\]

\(R_0\) solves a retained current equation; \(M_0\) solves a material normal equation. Put

\[
\begin{aligned}
\bar C&=(I-R_0L)C=C+R_0AC,& Y_C&=S\bar C,\\
N_C&=C^*(B-LK_0),& \Gamma_C&=C^*L\bar C.
\end{aligned}\tag{5}
\]

The symbol \(\Gamma_C\) denotes a current Schur matrix and is distinct from the receiver map \(S\). For \(A_{UV}=U^*AV\), \(D_V=I-A_{VV}\), \(B_V=V^*B\), and \(B_C=C^*B\), these factors become

\[
\begin{aligned}
\bar C&=C+VD_V^{-1}A_{VC},\\
Y_C&=SC+(SV)D_V^{-1}A_{VC},\\
N_C&=B_C+A_{CV}D_V^{-1}B_V,\\
\Gamma_C&=I-A_{CC}-A_{CV}D_V^{-1}A_{VC}.
\end{aligned} \tag{6}
\]

\textbf{Proposition 1 (feedback-dressed tangent factorization).} If \(\Gamma_C\) is invertible, the strictly Galerkin tangent in \([V,C]\) satisfies

\[
K_C-K_0=\bar C\Gamma_C^{-1}N_C,\qquad
\Delta J:=J_C-J_0=Y_C\Gamma_C^{-1}N_C. \tag{7}
\]

\textbf{Proof.} The enlarged coefficient matrix has blocks \(D_V,-A_{VC},-A_{CV},I-A_{CC}\). Eliminating its retained coordinates gives \(\Gamma_C\) and right-hand side \(N_C\). Back-substitution adds \(VD_V^{-1}A_{VC}\) to the raw candidate response \(C\), proving (6)--(7). It also gives \(V^*L\bar C=0\). \hfill\(\square\)

Equation (6) identifies three physical paths. First, an added current radiates directly through \(SC\), but also generates an internal field through \(G_D\), excites retained currents through \(X\), and radiates through \(SV\) after retained-space rescattering. Hence

\[
\boxed{\bar C=C+R_0XG_DC}
\]

in the contrast-source formulation. Second, material injection can enter the retained space through \(B_V\) before reaching the new coordinates through \(A_{CV}\). Third, \(A_{CV}D_V^{-1}A_{VC}\) is the added-to-retained-to-added loop, which enters the Schur inverse rather than a single-pass correction. The dressed current preserves the retained Galerkin equations through \(V^*L\bar C=0\), while generally losing Euclidean orthogonality to \(V\). For a receiver-dark candidate, the isolated direct branch \(SC\) is zero or near roundoff; internal scattering produces the nonzero dressed response \(S\bar C\).



\subsection{Weak relative feedback requires more than weak scattering}

If \(\|A_{VV}\|\le\eta<1\), a sufficient absolute bound is

\[
\|Y_C-SC\|\le
\frac{\|SV\|\,\|A_{VC}\|}{1-\eta}. \tag{8}
\]

To obtain a direction-uniform \emph{relative} bound on a candidate subspace, its direct radiation must also be bounded away from zero. If \(SC\) is injective there, then

\[
\frac{\|(Y_C-SC)a\|}{\|SCa\|}
\le\frac{\|SV\|\,\|A_{VC}\|}
{(1-\eta)\sigma_{\min}(SC)}
\quad(a\ne0). \tag{9}
\]

A singular \(SC\) requires restriction to a specified non-dark coefficient space. A sufficient Schur-invertibility condition is \(\|A_{CC}+A_{CV}D_V^{-1}A_{VC}\|<1\). These sufficient conditions separate weak absolute coupling from weak relative feedback. Small absolute feedback may dominate a nearly dark direct response; cross-space coupling and a large, possibly nonnormal retained resolvent provide additional amplification mechanisms. Their contributions must be resolved before attributing a large ratio to resonance.

\section{Paired Material Closure and Nonmonotone Curvature}

\subsection{The two-channel regularized step change}

For the material-space formulas below, the factors must be real-lifted together. Let $P$ be the number of illuminations, $\mathbf Y=I_P\otimes Y_C$, $\Gamma_{\rm blk}=I_P\otimes\Gamma_C$, and $\mathbf N=[N_{C,1}^\top,\ldots,N_{C,P}^\top]^\top$. Define $\mathcal R(F)=[\operatorname{Re}F,-\operatorname{Im}F;\operatorname{Im}F,\operatorname{Re}F]$ and $\widehat N=[\operatorname{Re}\mathbf N;\operatorname{Im}\mathbf N]$. Then $\widehat D=\mathcal R(\mathbf Y)\mathcal R(\Gamma_{\rm blk})^{-1}\widehat N$, up to the fixed stacking permutation. Throughout this material-space subsection, $Y_C,\Gamma_C,N_C$ denote these lifted factors, and $T_C=\Gamma_C^{-1}N_C$ is formed after lifting; their adjoints are real transposes. If $m$ is the complex receiver dimension per illumination and $r_c$ the added complex current rank, the lifted $Y_C$ has dimensions $2Pm\times2Pr_c$, and the lifted $N_C$ has dimensions $2Pr_c\times d$. This preserves the shared material variable and does not treat each complex current atom as a rank-one real material change.

Set \(T_C=\Gamma_C^{-1}N_C\). For the same \(B,r,\Lambda,\ell\), define

\[
\begin{aligned}
s_0&=-M_0(J_0^*r+\ell),\\
s_C&=-(J_C^*J_C+\Lambda)^{-1}(J_C^*r+\ell),\\
\mathcal E_0(C)&=\operatorname{Range}M_0[J_0^*Y_C,N_C^*].
\end{aligned}\tag{10}
\]

\textbf{Theorem 1 (paired step-change closure).} Under Proposition 1, with exact actions, positive regularization, and no nonzero signature removed,

\[
\begin{aligned}
s_C-s_0={}&-M_0J_0^*Y_C(T_Cs_C)\\
&-M_0N_C^*\Gamma_C^{-*}Y_C^*(r+J_Cs_C)
\in\mathcal E_0(C).
\end{aligned} \tag{11}
\]

\textbf{Proof.} Write \(\Delta H=J_C^*J_C-J_0^*J_0\) and \(\Delta h=\Delta J^*r\). Expanding \(J_C=J_0+Y_CT_C\) gives

\[
\Delta H=J_0^*Y_CT_C+T_C^*Y_C^*J_0+T_C^*Y_C^*Y_CT_C.
\]

Subtract the normal equations: \(H_0(s_C-s_0)=-\Delta Hs_C-\Delta h\). Collect the first observation term and the remaining \(T_C^*Y_C^*(r+J_Cs_C)\) term, and use \(T_C^*=N_C^*\Gamma_C^{-*}\). \hfill\(\square\)

The observation lift \(M_0J_0^*Y_C\) redistributes the update along old sensitivities when they couple to the added response. The injection lift \(M_0N_C^*\) pulls the updated data residual back through the new current excitation. Looking only at \(\Delta h\) retains the second channel and can miss the first channel's Hessian coupling.

The factors identify the available routes for material influence, while the coefficients in (11) determine how the present residual activates them. For example, when $r=0$ and $\ell=0$, both regularized steps vanish regardless of the derivative change. This separates structural coupling from its activation in a particular inverse step.

\subsection{Why one side is generally insufficient}

The exact illustration

\[
J_0=[1,0],\quad J_C=[1,1],\quad \Lambda=I,\quad r=1,\quad\ell=0
\]

gives \(s_0=(-1/2,0)^\top\), \(s_C=(-1/3,-1/3)^\top\), and \(s_C-s_0=(1/6,-1/3)^\top\). The observation lift is along the first material coordinate and the injection lift along the second. Neither alone contains the change. Although \(\Delta J=[0,1]\) has rank one,

\[
\Delta H=\begin{bmatrix}0&1\\1&1\end{bmatrix},\qquad\det\Delta H=-1. \tag{12}
\]

Both \(H_0=J_0^*J_0+\Lambda\) and \(H_C=J_C^*J_C+\Lambda\) remain positive definite under \(\Lambda\succ0\). The indefinite difference \(\Delta H\) describes redistribution of material curvature between these two well-posed regularized problems.

This algebraic illustration embeds in \(L=I\) Galerkin reduction: take $V=e_1$, $C=e_2$, $S=[1,1]$, and $B=I_2$. It isolates the least-squares effect: two-sided material coupling can occur even without scattering feedback. Feedback determines the electromagnetic factors; material least squares supplies the two-channel normal coupling.
